% ========================================================
%  《习近平新时代中国特色社会主义思想概论》  期末全真模拟试卷（含参考答案）
%  覆盖范围：导论、第一章至第十二章核心考点
%  题型：30单选(30分)+10多选(20分)+5简答(50分)
% ========================================================
\documentclass[12pt, a4paper]{article}

% ---- 中文支持 ----
\usepackage[UTF8]{ctex}

% ---- 数学 ----
\usepackage{amsmath, amssymb, amsthm}
\usepackage{mathtools}

% ---- 页面布局 ----
\usepackage[top=2.5cm, bottom=2.5cm, left=2.8cm, right=2.8cm]{geometry}

% ---- 表格 ----
\usepackage{booktabs, array, multirow, makecell}
\usepackage{tabularx}

% ---- 页眉页脚 ----
\usepackage{fancyhdr}
\usepackage{lastpage}
\pagestyle{fancy}
\fancyhf{}
\renewcommand{\headrulewidth}{0.6pt}
\lhead{\small 习近平新时代中国特色社会主义思想概论·期末全真模拟}
\rhead{\small 共 \pageref{LastPage} 页}
\cfoot{\small 第 \thepage\ 页}

% ---- 计数器与样式 ----
\usepackage{enumitem}
\usepackage{xcolor}
\usepackage{tcolorbox}
\tcbuselibrary{skins, breakable}

% ---- 填空线 ----
\newcommand{\blank}[1]{\underline{\hspace{#1}}}

% ---- 答案盒子 ----
\newtcolorbox{answerbox}[1][]{
  breakable,
  colback=gray!6,
  colframe=gray!50,
  fonttitle=\bfseries\small,
  title=参考答案,
  #1
}

% ---- 题号格式 ----
\newcounter{bigq}
\newcommand{\BigQ}[1]{%
  \stepcounter{bigq}%
  \vspace{6pt}%
  \noindent{\large\bfseries 第\chinese{bigq}大题\quad #1}%
  \vspace{4pt}\par
}

\newcounter{subq}[bigq]
\newcounter{smallq}[subq]

\newcommand{\SmallQ}{%
  \stepcounter{smallq}%
  \noindent\arabic{smallq}.\quad
}

\newcommand{\resetsmallq}{\setcounter{smallq}{0}}

% 分隔线
\newcommand{\examrule}{\noindent\rule{\linewidth}{0.6pt}}

% 选项排版
\newcommand{\choices}[4]{\par
\noindent\hspace*{2em}A. #1\quad B. #2\quad C. #3\quad D. #4\par}

\newcommand{\choicelong}[4]{\par
\noindent\hspace*{2em}A. #1\par\hspace*{2em}B. #2\par\hspace*{2em}C. #3\par\hspace*{2em}D. #4\par}

\newcommand{\choicelongfive}[5]{\par
\noindent\hspace*{2em}A. #1\par\hspace*{2em}B. #2\par\hspace*{2em}C. #3\par\hspace*{2em}D. #4\par\hspace*{2em}E. #5\par}

% ---- 文档开始 ----
\begin{document}

% ============================================================
%  封面
% ============================================================
\begin{center}
  {\LARGE\bfseries 习近平新时代中国特色社会主义思想概论}\\[8pt]
  {\large\bfseries 期末全真模拟试卷·GLM·A卷}\\[16pt]
  \begin{tabular}{lll}
    考试时间：120 分钟 & \quad & 满分：100 分 \\[4pt]
    \multicolumn{3}{l}{注意事项：请将答案写在答题纸指定区域，写在本卷上无效。}
  \end{tabular}
  \vspace{4pt}

  \examrule

  \vspace{6pt}
  \begin{tabular}{|c|c|c|c|c|}
    \hline
    \textbf{题号} & 一 & 二 & 三 & \textbf{总分} \\
    \hline
    \textbf{满分} & 30 & 20 & 50 & 100 \\
    \hline
    \textbf{得分} & & & & \\
    \hline
  \end{tabular}
\end{center}

\vspace{10pt}
\examrule
\vspace{6pt}

% ============================================================
%  第一大题：单项选择题（30题，每题1分，共30分）
% ============================================================
\BigQ{单项选择题（每题 1 分，共 30 分。每题只有一个正确选项）}

\resetsmallq
\SmallQ 习近平新时代中国特色社会主义思想创立的时代背景之一是（\quad）。
\choices{世界百年未有之大变局加速演进}{冷战格局加剧}{工业革命兴起}{经济全球化逆转}
\vspace{4pt}
\SmallQ "两个确立"是指（\quad）。
\choicelong{确立习近平同志党中央的核心、全党的核心地位；确立习近平新时代中国特色社会主义思想的指导地位}{确立党的基本路线；确立党的基本方略}{确立中国特色社会主义道路；确立理论自信}{确立"两个维护"；确立"两个结合"}
\vspace{4pt}
\SmallQ 中国特色社会主义进入新时代，这是我国发展新的（\quad）。
\choices{发展阶段}{历史方位}{历史阶段}{历史时期}
\vspace{4pt}
\SmallQ 中国最大的国情是（\quad）。
\choices{中国共产党的领导}{人口规模巨大}{处于社会主义初级阶段}{发展中大国}
\vspace{4pt}
\SmallQ 中国共产党自身优势中，马克思主义政党第一位的能力是（\quad）。
\choices{政治领导力}{思想引领力}{群众组织力}{社会号召力}
\vspace{4pt}
\SmallQ 人民立场是中国共产党的（\quad）。
\choices{基本工作方法}{重要思想来源}{根本政治立场}{组织原则}
\vspace{4pt}
\SmallQ 共同富裕是中国特色社会主义的（\quad）。
\choices{本质要求}{基本制度}{现阶段任务}{长远规划}
\vspace{4pt}
\SmallQ 全面深化改革的总目标是（\quad）。
\choicelong{建立社会主义市场经济体制}{实现全面建成小康社会}{构建新发展格局}{完善和发展中国特色社会主义制度，推进国家治理体系和治理能力现代化}
\vspace{4pt}
\SmallQ 新发展理念中，注重解决发展动力问题的是（\quad）。
\choices{创新}{协调}{绿色}{开放}
\vspace{4pt}
\SmallQ "科技是第一生产力、人才是第一资源、创新是第一动力"的论断体现了（\quad）。
\choices{经济体制改革的核心}{教育、科技、人才一体推进的战略思想}{文化建设的根本任务}{党的建设总要求}
\vspace{4pt}
\SmallQ 人民民主是社会主义的（\quad）。
\choices{生命}{本质}{特征}{核心}
\vspace{4pt}
\SmallQ 下列表述属于"两个结合"内容的是（\quad）。
\choicelong{马克思主义基本原理同中国具体实际相结合、同中华优秀传统文化相结合}{马克思主义基本原理同中国革命相结合、同建设相结合}{马克思主义基本原理同新时代相结合、同改革开放相结合}{马克思主义基本原理同东方文化相结合、同世界潮流相结合}
\vspace{4pt}
\SmallQ 习近平新时代中国特色社会主义思想是"两个结合"的重大成果，其中"魂脉"是（\quad）。
\choices{中华优秀传统文化}{马克思主义基本原理}{中国具体实际}{社会主义先进文化}
\vspace{4pt}
\SmallQ 在习近平新时代中国特色社会主义思想科学体系中，起支撑作用的"四梁八柱"是（\quad）。
\choices{十个明确}{十四个坚持}{十三个方面成就}{六个必须坚持}
\vspace{4pt}
\SmallQ "十四个坚持"是新时代坚持和发展中国特色社会主义的（\quad）。
\choices{根本制度}{基本路线}{基本方略}{战略目标}
\vspace{4pt}
\SmallQ 道路问题是关系党的事业兴衰成败的（\quad）的问题。
\choices{第一位}{核心}{关键}{基础}
\vspace{4pt}
\SmallQ "一个变与两个没有变"中，发生变化的是（\quad）。
\choices{基本国情}{社会主要矛盾}{国际地位}{社会性质}
\vspace{4pt}
\SmallQ 我国仍是世界最大发展中国家的国际地位（\quad）。
\choices{没有变}{已经改变}{即将改变}{部分改变}
\vspace{4pt}
\SmallQ 党的基本路线是国家的（\quad）。
\choices{生命线、人民的幸福线}{根本制度}{基本方略}{行动指南}
\vspace{4pt}
\SmallQ 中国梦的本质是（\quad）。
\choices{全面建设小康社会}{实现现代化}{中华民族伟大复兴}{国家富强、民族振兴、人民幸福}
\vspace{4pt}
\SmallQ "五位一体"总体布局中不包括（\quad）。
\choices{党的建设}{经济建设}{文化建设}{生态文明建设}
\vspace{4pt}
\SmallQ "五个必由之路"中，团结奋斗是中国人民创造历史伟业的（\quad）。
\choices{必由之路}{首要任务}{根本制度}{战略目标}
\vspace{4pt}
\SmallQ 民心是最大的政治，决定事业（\quad）。
\choices{兴衰成败}{方向道路}{速度规模}{理论根基}
\vspace{4pt}
\SmallQ 中国共产党始终坚守的根本工作方法是（\quad）。
\choices{顶层设计}{群众路线}{试点先行}{试点推广}
\vspace{4pt}
\SmallQ 社会主义基本经济制度新概括不包括（\quad）。
\choices{计划经济体制}{公有制为主体、多种所有制经济共同发展}{按劳分配为主体、多种分配方式并存}{社会主义市场经济体制}
\vspace{4pt}
\SmallQ 新发展格局是（\quad）。
\choices{以国内大循环为主体、国内国际双循环相互促进}{以国际大循环为主体}{封闭的国内单循环}{完全开放型经济}
\vspace{4pt}
\SmallQ 全过程人民民主的"四个相统一"不包括（\quad）。
\choices{过程民主和成果民主的统一}{过程民主和结果民主的统一}{程序民主和实质民主的统一}{人民民主和国家意志的统一}
\vspace{4pt}
\SmallQ 中国特色社会主义文化的三个主要组成部分不包括（\quad）。
\choices{中华优秀传统文化}{革命文化}{资本主义文化}{社会主义先进文化}
\vspace{4pt}
\SmallQ 意识形态关乎旗帜、关乎道路、关乎（\quad）。
\choices{国家安全}{经济发展}{社会稳定}{民生福祉}
\vspace{4pt}
\SmallQ 生态文明建设应摆在（\quad）。
\choices{经济工作的次要位置}{文化建设的基础位置}{社会建设的边缘位置}{全局工作的突出位置}

\vspace{10pt}
\examrule

% ============================================================
%  第二大题：多项选择题（10题，每题2分，共20分）
% ============================================================
\BigQ{多项选择题（每题 2 分，共 20 分。每题有两个或两个以上正确选项，少选、错选均不得分）}

\resetsmallq
\SmallQ 习近平新时代中国特色社会主义思想创立的时代背景包括（\quad）。
\choicelong{世界百年未有之大变局加速演进}{中华民族伟大复兴进入关键时期}{中国式现代化全面推进拓展}{科学社会主义在21世纪的中国焕发新的蓬勃生机}
\vspace{4pt}
\SmallQ 习近平新时代中国特色社会主义思想科学体系包括（\quad）。
\choices{十个明确}{十四个坚持}{十三个方面成就}{六个必须坚持}
\vspace{4pt}
\SmallQ 推进中国式现代化必须牢牢把握的重大原则包括（\quad）。
\choicelong{坚持和加强党的全面领导}{坚持中国特色社会主义道路}{坚持以人民为中心的发展思想}{坚持深化改革开放}
\vspace{4pt}
\SmallQ 党的自身优势包括（\quad）。
\choices{政治领导力}{思想引领力}{群众组织力}{社会号召力}
\vspace{4pt}
\SmallQ 新发展理念的科学内涵包括（\quad）。
\choices{创新}{协调}{绿色}{开放}
\vspace{4pt}
\SmallQ 全过程人民民主的"四个相统一"包括（\quad）。
\choicelong{过程民主和成果民主的统一}{程序民主和实质民主的统一}{直接民主和间接民主的统一}{人民民主和国家意志的统一}
\vspace{4pt}
\SmallQ 促进共同富裕要把握好的原则包括（\quad）。
\choices{鼓励勤劳创新致富}{坚持基本经济制度}{尽力而为量力而行}{坚持循序渐进}
\vspace{4pt}
\SmallQ 全面深化改革开放要坚持正确的方法论，包括（\quad）。
\choicelong{增强系统性、整体性、协同性}{加强顶层设计和摸着石头过河相结合}{统筹改革发展稳定}{坚持重大改革于法有据}
\vspace{4pt}
\SmallQ 中国式现代化创造了人类文明新形态，主要表现在（\quad）。
\choicelong{提供全新现代化路径，超越西方现代化模式}{为发展中国家提供新的选择}{需牢牢把握推进中国式现代化的重大原则}{完全照搬西方现代化道路}
\vspace{4pt}
\SmallQ 关于"一个变与两个没有变"的表述，下列正确的是（\quad）。
\choicelong{社会主要矛盾已经变化}{我国仍处于社会主义初级阶段的基本国情没有变}{我国仍是世界最大发展中国家的国际地位没有变}{我国社会性质已经发生根本变化}

\vspace{10pt}
\examrule

% ============================================================
%  第三大题：简答题（5题，每题10分，共50分）
% ============================================================
\BigQ{简答题（每题 10 分，共 50 分。要求观点正确、要点完整、简要论述）}

\resetsmallq
\SmallQ（10 分）简述习近平新时代中国特色社会主义思想创立的时代背景。
\vspace{6pt}

\SmallQ（10 分）简述"两个确立"的决定性意义。
\vspace{6pt}

\SmallQ（10 分）简述推进中国式现代化必须牢牢把握的重大原则。
\vspace{6pt}

\SmallQ（10 分）为什么说新时代全面深化改革开放是一场深刻革命？
\vspace{6pt}

\SmallQ（10 分）简述全过程人民民主的"四个相统一"。

\vspace{10pt}
\examrule
\vspace{16pt}
\examrule
\vspace{4pt}
\begin{center}
  {\large\bfseries ——试题到此结束，请认真检查——}
\end{center}
\vspace{4pt}
\examrule

% ============================================================
%  参考答案
% ============================================================
\newpage
\setcounter{bigq}{0}

\begin{center}
  {\LARGE\bfseries 参考答案与评分标准}
\end{center}
\vspace{8pt}
\examrule

% -------- 第一大题参考答案 --------
\BigQ{单项选择题参考答案}

\begin{answerbox}

\begin{tabular}{@{}llllll@{}}
1. A & 2. A & 3. B & 4. A & 5. A & 6. C \\
7. A & 8. D & 9. A & 10. B & 11. A & 12. A \\
13. B & 14. A & 15. C & 16. A & 17. B & 18. A \\
19. A & 20. D & 21. A & 22. A & 23. A & 24. B \\
25. A & 26. A & 27. B & 28. C & 29. A & 30. D \\
\end{tabular}

\vspace{6pt}
\textbf{要点解析（节选）：}

\begin{itemize}[leftmargin=2em,itemsep=2pt]
  \item 第1题：时代背景"五个方面"第一条即"世界百年未有之大变局加速演进"。
  \item 第2题："两个确立"是党确立习近平同志党中央的核心、全党的核心地位，确立习近平新时代中国特色社会主义思想的指导地位。
  \item 第3题：新时代是我国发展新的"历史方位"，注意与"历史阶段"等表述区分。
  \item 第4题：中国最大的国情就是中国共产党的领导，这是第三章第一考点。
  \item 第5题：政治领导力是马克思主义政党第一位的能力。
  \item 第6题：人民立场是中国共产党的根本政治立场，是党的显著标志。
  \item 第7题：共同富裕是中国特色社会主义的本质要求，是中国式现代化的重要特征。
  \item 第8题：总目标由"完善和发展中国特色社会主义制度"与"推进国家治理体系和治理能力现代化"两句话构成。
  \item 第9题：创新是引领发展的第一动力，注重解决发展动力问题。
  \item 第10题："三个第一"体现了教育、科技、人才一体推进的战略思想。
  \item 第11题："人民民主是社会主义的生命"是第八章第一考点。
  \item 第12题："两个结合"是马克思主义基本原理同中国具体实际相结合、同中华优秀传统文化相结合。
  \item 第13题：马克思主义基本原理是"魂脉"，中华优秀传统文化是"根脉"，注意区分。
  \item 第14题："十个明确"是这一思想最为核心关键的组成部分，是"四梁八柱"。
  \item 第15题："十四个坚持"也称"十四个方略"，是新时代坚持和发展中国特色社会主义的基本方略。
  \item 第16题：道路问题是关系党的事业兴衰成败的"第一位"的问题。
  \item 第17题："一个变"是社会主要矛盾的变化；"两个没有变"是基本国情和国际地位。
  \item 第18题：我国仍是世界最大发展中国家的国际地位"没有变"。
  \item 第19题：党的基本路线是国家的生命线、人民的幸福线。
  \item 第20题：中国梦的本质是国家富强、民族振兴、人民幸福。
  \item 第21题："五位一体"包括经济、政治、文化、社会、生态文明建设，不包括党的建设。
  \item 第22题：团结奋斗是中国人民创造历史伟业的必由之路，是"五个必由之路"之一。
  \item 第23题：民心是最大的政治，决定事业兴衰成败。
  \item 第24题：群众路线是党始终坚守的根本工作方法；调查研究是重要工作方法。
  \item 第25题：基本经济制度新概括由所有制、分配制度、社会主义市场经济体制构成，计划经济体制不属于。
  \item 第26题：新发展格局以国内大循环为主体、国内国际双循环相互促进。
  \item 第27题：四个相统一是过程民主与成果民主、程序民主与实质民主、直接民主与间接民主、人民民主与国家意志；不存在"结果民主"。
  \item 第28题：文化三部分为中华优秀传统文化、革命文化、社会主义先进文化。
  \item 第29题：意识形态关乎旗帜、关乎道路、关乎国家安全。
  \item 第30题：生态文明建设应摆在全局工作的突出位置。
\end{itemize}

\end{answerbox}

% -------- 第二大题参考答案 --------
\BigQ{多项选择题参考答案}

\begin{answerbox}

\begin{tabular}{@{}ll@{}}
1. ABCD & 2. ABCD \\
3. ABCD & 4. ABCD \\
5. ABCD & 6. ABCD \\
7. ABCD & 8. ABCD \\
9. ABC & 10. ABC \\
\end{tabular}

\vspace{6pt}
\textbf{要点解析：}

\begin{itemize}[leftmargin=2em,itemsep=2pt]
  \item 第1题：时代背景五个方面中，本题四个选项均为正确表述。
  \item 第2题：思想体系由"十个明确、十四个坚持、十三个方面成就、六个必须坚持"构成。
  \item 第3题：五大重大原则全部列出，全选。
  \item 第4题：党的自身优势即"政治领导力、思想引领力、群众组织力、社会号召力"四力。
  \item 第5题：新发展理念五字诀为"创新、协调、绿色、开放、共享"，本题列出创新、协调、绿色、开放，全选。
  \item 第6题：四个相统一全部为正确表述。
  \item 第7题：共同富裕四原则——鼓励勤劳创新致富、坚持基本经济制度、尽力而为量力而行、坚持循序渐进。
  \item 第8题：方法论五条全部为正确表述。
  \item 第9题：人类文明新形态三方面表现——提供全新现代化路径超越西方、为发展中国家提供新选择、需把握五大重大原则；D项"照搬西方"错误，故选ABC。
  \item 第10题："一个变"是社会主要矛盾的变化；"两个没有变"是基本国情和国际地位。社会性质未变，故D项错误，选ABC。
\end{itemize}

\end{answerbox}

% -------- 第三大题参考答案 --------
\BigQ{简答题参考答案}

\begin{answerbox}

\textbf{第1题（10分）：习近平新时代中国特色社会主义思想创立的时代背景}

习近平新时代中国特色社会主义思想创立的时代背景主要包括以下五个方面（每点2分，共10分）：

\textbf{（1）}世界百年未有之大变局加速演进。国际力量对比发生深刻变化，全球治理体系与国际秩序深度调整，新一轮科技革命和产业变革深入发展，世界进入新的动荡变革期。

\textbf{（2）}中华民族伟大复兴进入关键时期。我国发展站在新的历史起点上，面临前所未有的机遇和挑战，需要新的思想引领伟大事业继续前进。

\textbf{（3）}中国式现代化全面推进拓展。党领导人民成功推进和拓展了中国式现代化，开辟了人类现代化的新道路。

\textbf{（4）}科学社会主义在21世纪的中国焕发新的蓬勃生机。中国特色社会主义使科学社会主义在21世纪的中国焕发新的蓬勃生机，展现了强大生命力。

\textbf{（5）}中国共产党在革命性锻造中更加坚强有力。党在新时代的自我革命中更加坚强有力，为创立新思想提供了坚强的政治保证。

\end{answerbox}

\begin{answerbox}

\textbf{第2题（10分）："两个确立"的决定性意义}

"两个确立"是指确立习近平同志党中央的核心、全党的核心地位，确立习近平新时代中国特色社会主义思想的指导地位。（2分）

其决定性意义主要体现在：

\textbf{（1）}"两个确立"是新时代取得一切成就的根本原因。党的十八大以来党和国家事业取得历史性成就、发生历史性变革，最根本的原因在于有习近平总书记作为党中央的核心、全党的核心掌舵领航，在于有习近平新时代中国特色社会主义思想科学指引。（3分）

\textbf{（2）}"两个确立"是战胜一切艰难险阻、应对一切不确定性的最大确定性。面对复杂多变的国际国内形势，"两个确立"为党和国家事业发展提供了根本政治保证。（2分）

\textbf{（3）}"两个确立"是最大底气、最大保证。它体现了全党全军全国各族人民共同心愿，对新时代党和国家事业发展、对推进中华民族伟大复兴历史进程具有决定性意义。（3分）

\end{answerbox}

\begin{answerbox}

\textbf{第3题（10分）：推进中国式现代化必须牢牢把握的重大原则}

推进中国式现代化必须牢牢把握以下五个重大原则（每点2分，共10分）：

\textbf{（1）}坚持和加强党的全面领导。这是中国式现代化的根本保证，要把党的领导贯彻到推进中国式现代化的全过程各方面。

\textbf{（2）}坚持中国特色社会主义道路。这是中国式现代化的根本方向，要始终沿着中国特色社会主义道路前进，既不走封闭僵化的老路，也不走改旗易帜的邪路。

\textbf{（3）}坚持以人民为中心的发展思想。这是中国式现代化的根本立场，要使中国式现代化的成果更多更公平惠及全体人民。

\textbf{（4）}坚持深化改革开放。这是中国式现代化的根本动力，要以改革开放作为推动中国式现代化的强大引擎。

\textbf{（5）}坚持发扬斗争精神。这是中国式现代化的根本要求，要增强全党全国各族人民的志气、骨气、底气，依靠顽强斗争打开事业发展新天地。

\end{answerbox}

\begin{answerbox}

\textbf{第4题（10分）：为什么说新时代全面深化改革开放是一场深刻革命？}

新时代全面深化改革开放是一场深刻革命，主要体现在以下几个方面：

\textbf{（1）}从其艰巨性、复杂性和系统性来说，是一场深刻的革命。新时代改革开放涉及面广、触及利益深、攻坚难度大，其艰巨性、复杂性、系统性前所未有。（2分）

\textbf{（2）}全面深化改革开放是涉险滩、闯难关的变革。它要破解体制机制弊端、突破利益固化藩篱，触及深层次利益格局调整，是一场需要啃硬骨头的革命。（2分）

\textbf{（3）}全面深化改革开放是一场思想理论的深刻变革、改革组织方式的深刻变革、国家制度和治理体系的深刻变革、人民广泛参与的深刻变革。（3分）

\textbf{（4）}全面深化改革开放是在多年改革开放基础上的深化，是一场全面、系统、整体的制度创新。它要求着眼改革的关联性、系统性、可行性，全面系统整体推进各方面改革，使社会主义制度的优越性得到更好发挥。（2分）

\textbf{（5）}新时代的改革开放是全方位、深层次、根本性的，取得的成就是历史性、革命性、开创性的，深刻改变了中国的面貌，也深刻影响了世界。（1分）

\end{answerbox}

\begin{answerbox}

\textbf{第5题（10分）：全过程人民民主的"四个相统一"}

全过程人民民主是社会主义民主政治的本质属性，其"四个相统一"具体如下（每点2.5分，共10分）：

\textbf{（1）}全过程人民民主是过程民主和成果民主的统一。它不仅关注民主程序的完整性，更注重民主成果的实效性，既重视过程参与，又重视结果落实，体现了过程与成果的辩证统一。

\textbf{（2）}全过程人民民主是程序民主和实质民主的统一。它既通过完整的制度程序保障民主运转，又通过实质性的民主内容体现人民意志，使民主不流于形式，而是真正落到实处。

\textbf{（3）}全过程人民民主是直接民主和间接民主的统一。它通过人民代表大会制度等间接民主形式和基层群众自治等直接民主形式相结合，形成完整的民主参与体系。

\textbf{（4）}全过程人民民主是人民民主和国家意志的统一。它将人民意愿通过法定程序转化为国家意志，使人民当家作主与国家治理有机统一，体现了社会主义民主政治的优越性。

\end{answerbox}

\vspace{12pt}
\examrule
\begin{center}
  \textit{参考答案仅供参考，评分时应酌情给分，步骤正确可得过程分。}

\end{center}

\end{document}